%        File: cv.tex
%     Created: Fri Sep 21 10:00 AM 2012 C
% Last Change: Fri Sep 21 10:00 AM 2012 C
%
%
%
%  CRAZY


% Text dimensions



\documentclass{article}
\usepackage{newspaper}

\date{\today}
\currentvolume{1}
\currentissue{73}

\usepackage{times}
\usepackage[dvipdfmx]{graphicx}
\usepackage{multicol}
\usepackage{picinpar}

%\usepackage{amsmath,amsfonts,amsthm} 
\usepackage[svgnames]{xcolor} 
%\usepackage{geometry} 
%%     \textheight=700px 
%\usepackage{url} 
\usepackage{wrapfig}
%
%\frenchspacing
%\pagestyle{empty}
%
\usepackage{sectsty}
  \sectionfont{%
    \usefont{OT1}{phv}{b}{n}%
    \sectionrule{0pt}{0pt}{-5pt}{3pt}}
%
\begin{document}
\newlength{\spacebox}
\settowidth{\spacebox}{1234567890}
%
\newcommand{\sepspace}{\vspace*{1em}}
%
\newcommand{\MyName}[1]{
  \Huge \usefont{OT1}{phv}{b}{n} \hfill #1
  \par \normalsize \normalfont
}

\newcommand{\MySlogan}[1]{
  \small \usefont{OT1}{phv}{m}{n}\hfill \textit{#1} % Slogan (optional)
  \par \normalsize \normalfont
}

\newcommand{\NewPart}[1]{\section*{\uppercase{#1}}}
%
\newcommand{\PersonalEntry}[2]{
  \noindent\hangindent=2em\hangafter=0
  \parbox{\spacebox}{
    \textit{#1}
  }
  \hspace{1.5em} #2 \par
}
%
\newcommand{\EducationEntry}[4]{
		\noindent \textbf{#1} \hfill 					% Study
		\colorbox{Black}{%
			\parbox{6em}{%
			\hfill\color{White}#2}} \par				% Duration
		\noindent \textit{#3} \par				    	% School
		\noindent\hangindent=2em\hangafter=0 \small #4 	% Description
		\normalsize \par}
%
\newcommand{\WorkEntry}[4]{						    % Same as \EducationEntry
		\noindent \textbf{#1} \hfill 				% Jobname
		\colorbox{Black}{\color{White}#2} \par		% Duration
		\noindent \textit{#3} \par					% Company
		\noindent\hangindent=2em\hangafter=0 \small #4 	% Description
		\normalsize \par}
%
%%\setlength\topmargin{-48pt}
%%\setlength\headheight{0pt}
%%\setlength\headsep{25pt}
%%\setlength\marginparwidth{-20pt}
%%\setlength\textwidth{7.0in}
%%\setlength\textheight{9.5in}
%%\setlength\oddsidemargin{-30pt}
%%\setlength\evensidemargin{-30pt}
%
\SetPaperName{Mikael TimeS}
\SetHeaderName{Awesome}
\SetPaperLocation{Malm\"{o}}
%\SetPaperSlogan{Working to create}
\SetPaperPrice{USD $10^6$}
%
%
%\author{Mikael Vatau}
%\title{CV}
\maketitle
\begin{multicols}{3}{
    \byline{\sc\Large The Essence of Tools}{Mikael Vatau}

    \begin{window}[2,r,\includegraphics[width=1in]{mesmall.eps},]
    There is a lot to be said about using the right tools for the job, says
    Mikael, a software engineer who has worked for over 12 years with embedded
    systems. Chefs brings their own favourad kitchen tools to work, it should
    be the same for programmers. Mikael worked as a software developer in the 
    mobile phone space for 7 years and before that as a research engineer at 
    robotics department at Lunds University.

    When asked what his favoured tools are he answers vim, zsh, linux, brain,
    coffe, creativity and optimism, and ofcourse adaptation. \textit{''You 
    think about what you are trying to solve and then how to get there.
    And then you choose your tools.''} 
    \end{window}
  \closearticle

  \headline{\it\huge Lund University}
    Mikael has worked in both small and large organisations. When working at the
    Robotics department at Lund University, he was the only expert on embedded
    system. This was the first job after University  and there was alot of
    responsibility. He had to design the schematics for the embedded systems
    and then produce the software for these. \textit{''The differense between
    designing hardware and software was the development cycles. You get almost
    instant response if your program works but you had to wait for  weeks to get
    feedback that your hardware design worked''}.
    
    Another aspect he remembered from this time was that he missed having 
    colleges working in the same area to share solution with. But that was also
    the good part. \textit{''I learned alot and I mostly made the right decision, a 
    confident booster. I was also involved in a lot of projects and met many
    intresting people during my time there.''} 
    
    After that period he started to work as a consultant and mostly worked in 
    larger mobile phone companies. \textit{''I started out as a integretor at 
    Ericsson AB. My boss told me that my work was being a cop. I integrated 
    other peoples software and when someone broke the platform you hunted down
    the guilty one and made him fix it or remove his part. The debugging part
    and was fun, but documenting what went in was very monotones work.''} 

    He also was responsible for the lint tool, deciding what to filter, what 
    to allow and answering questions from other developers about the tools
    warning. \textit{''That was fun and educational. I learned alot about good
    program practices and the quirks in the C language.''}

    But he was glad when he startet to work as a software developer again. 
    \textit{``Programming is enjoyable. Spending hours, days, weeks, to 
    improve the code. Cleaner more efficient code.  Shaving of a few 
    milliseconds, use less resources, making it more perfect, more beatifull.
    Having spend that time, editing some text according to specific rules. 
    Your code is fluctual, might be thrown out the next rewrite,
    might stay there a long time. Might be ignored. Might be used by 
    millions. Sorry, getting to filosofical? ''} 

    He started on a new assignment in 2009 to Sony Ericsson, where he worked
    with the ETS software used to test the phones at the factory. This was 
    mostly low level development using the android development. \textit{''We got
    to play with the new hardware, often helping out to get the drivers work 
    with the hardware. I was mosltly involved with the GPS tests, Since 
    qualcomms testsoftware did not include a way to calculate the S/N ratio for
    the GPS signal on the phone and we used the serial port to extract the S/N
    ratio  this test took several minutes to performe, so I ported the S/N ratio
    calculation to the phone and got down the test to 4 seconds. There was
    optimisation involved which made it enjoyable. I also made some of the
    display and touch tests, and helped a little with the wlan driver''}

    Then he started working at Ericsson again, in audio this time, working with
    Low Power Audio the alsa driver, fixing error reports and traveling to
    customer for support. \textit{''The thing I will remember at this job is 
    traveling the work with the customers, the differens in working culture. 
    It was a tight schedule and lots of bugs. Somehow we made it work though. We
    had good support team back home and my most important role was not to solve 
    bugs, but to work with customer relationship. It was exciting time with lots
    of action and a good outcome. Coming home to the slow scandinavian pace I
    almost got post war syndrome, haha.''}

    Now he is working with a new project, where Sony Mobile is outsourcing
    several Android applications. One of the places is to the consulting firm
    where he works. 

}

\end{multicols}

%\begin{wrapfigure}{l}{0.5\textwidth}
%	\vspace*{-2em}
%		\includegraphics[width=0.15\textwidth]{mesmall.eps}
%\end{wrapfigure}

\MyName{Mikael Vatau}
\MySlogan{work creates, creativity is fun} %mv why doesn't it say what I want damn it..,

\sepspace

\NewPart{Me} %mv should probably at some skills here..
\PersonalEntry{Past work}{Embedded systems, Mobile phones, Web, Robotics}
\PersonalEntry{Languages}{C, C++, Java, PHP, Pascal, English, Swedish}
\PersonalEntry{Tools}{zsh, vim, internet and coffee. Otherwise what will fit best}
\PersonalEntry{Likes}{Thinking, Optimizing, Positivity, Warmth, Good food}

\NewPart{Netville, consultant (2005 - present)}{Consulting for more than 
  7 years for different customers in the mobile phone industry. Worked in many 
  different areas such as software configuration, developer for low level
  audio, GPS, touch screen, WLAN, Bluetooth, as well as on-site support and
higher level application development.}
\sepspace

\WorkEntry{Developer}{Oct 2010-present}{ST-Ericsson}
  {Helped debug and solve audio issues working with the Low level
    audio-hardware team on the Android platform at the kernel level close to
    the hardware. I was involved in Low power audio playback and the wideband
    Bluetooth usecases and also worked with the LTP tests for the ALSA driver. 
  I supported Ericssons customers in both Asia and Sweden on site with audio
  issues before releases. 
  the ALSA driver.}
\sepspace

\WorkEntry{Developer}{Jun 2010-Sep 2010}{Netville Customer}
{Developed an Android application to help visually impaired people making
  phone calls.} 
\sepspace

\WorkEntry{Developer}{Mar 2009 - Apr 2010}{Sony Ericsson Mobile Communications
            AB}
  {Developer of Embedded Test Software (ETS) for testing hardware components in
    mobile phones and worked mainly with GPS, WLAN and touch displays. The ETS
    software was used at the production site both when manufacturing developer
    prototypes and in the high volume manufacturing. The ETS software was
    running on an Android platform and since we were one of the first to get
    involved with new platforms, we also helped with bringing up the drivers.}
\sepspace

\WorkEntry{Developer}{Mar 2005 - May 2008}{Ericsson AB}	
  {Supported developers with the PCLint tool. is a static code analysis tool
    that flag suspicious constructs in the c code and was used by all the
    developers before releasing code changes. I worked with configuring the
    PCLint tool for new platform versions (OSE) and when changing compilers. I
    also supported developers with advice concerning ANSI C and/or updating the
    configuration files. 
  }
\sepspace

\WorkEntry{Developer}{Mar 2005 - May 2007}{Ericsson AB}		
  {Worked with integration of component subsystems in a ClearCase environment
    at the Open Platform Architecture (OPA) department. This included
    regression testing, release and test documentation.  When there was a
    regression I was responsible of finding out the guilty party by performing
    log-, dump analysis analysis and debugging with Lauterbach}
\sepspace

\NewPart{Robotics, Lund University (2000 - 2004)}{Worked at the Robotics
  department, Lund University, in almost 5 years with several project, I had
  the responsibility for the hardware and software and was part of the team
  that administrated the departments computer and developed the department new
  website.}
\sepspace

\WorkEntry{System administration}{Jan 2000 - Sep 2004}{Robotics institution,
  University of Lund}		
  {Administrating IRIX, Linux and Windows Network at the robotics department.}
\sepspace

\WorkEntry{Developer}{Apr 2003 - Sep 2004}{Robotics institution, University of Lund}		
  {Developed software for a laboratory robot used by AstraZeneca for protein testing. 
  The robot would take a protein sample, deploy it, and then clean the picker. 
  Software was done in PLC using the IsaGRAPH Designer.}
\sepspace

\WorkEntry{Developer}{Jul 2003 - Feb 2004}{Robotics institution, University of Lund}		
  {Developed a new website for the Robotics department which included news
    section, document repository.  PHP, MySQL and JavaScript.} 
\sepspace

\WorkEntry{HW design}{Jan 2000 - Jun 2003}{Robotics institution, University of Lund}		
{Designed several variants of micro controller cards that where used to control
  DC and Step motors. The microcontroller card was based on a 32 bit CPU from
  Freescale. The design was made in Protel Cad software and verify /
  troubleshoot the hardware. Different versions were produced depending
  on power, size, and communication requirements.
}
\sepspace

\WorkEntry{Developer}{Mar 2000 - Jun 2003}{Robotics institution, University of Lund}		
{Developed the communication software used on an embedded system to communicate
  both with other  embedded systems and with a PC via the CAN protocol.
  The software was for a microcontroller that would be used to control the
  nodes on a robotic arm, attached to a wheelchair for physically impaired
  people} 
\sepspace

\WorkEntry{Developer}{Jan 2001 - Jun 2002}{Robotics institution, University of Lund}		
{Responsible for the electronics and software in a EU project where we
  developed a can/bottle opener that would open the container at a push of a
  button. The can opener was for impaired people so safety was also involved.}


\NewPart{Education}{}

\EducationEntry {Linux for Embedded Systems Advanced}{2009}{Sony Ericsson
  Mobile Communications AB (ENEA)}
{}	
\sepspace

\EducationEntry {Linux for Embedded Systems}{2008}{Ericsson AB}
{}	
\sepspace

\EducationEntry {GPRS Security}{2006}{Ericsson AB}
{}	
\sepspace

\EducationEntry{Practical English in Cross-cultural Communications}{2003}
{malm\"{o}ö University}
{}
\sepspace

\EducationEntry {IRIX System administration}{2000}{SGI}
{}	
\sepspace

\EducationEntry{B.Sc.E.E. Electro Engineer}{1997-2000}
{Lund Institute of Technology} {}

\end{document}


